% 6.电路基础.tex 子文件，book文档类，主文档通过\input{6.电路基础.tex}导入
\chapter{第五章 三相电路}

\section{三相电源和三相电路}
如果三相电压不仅频率相同、波形和幅值相同，而且变化进程的时间差相等，即彼此相差，则称为对称三相电压。

如果三相电压，其瞬时表达式为
\[
\]
取$k=1$，则为
\[
\]
其波形如图所示，对称正弦三相电压有效值相量为
\[
\]
式中，是一个旋转因子。

三相电压的相位。式中三相电压的相位的次序为，称为正序或顺序。

由三相电源、三相负载和连接导线所组成的电路称为三相电路。

\section{星形联结和三角形联结}
\subsection{星形联结}
图中三相电源和负载均为星形，或称$Y$形。三个相连在一起的公共点称为中性点。中性点之间的连线称为中性线。在这种连接方式中，电源。

在对称三相电路中，各相电流必然是对称的。此时
\[
\]
上式说明。

端线电流。

根据星形联结的电路结构可知：后者用相量表示为
\[
\]

在相电压对称的情况下，线电压与相电压的关系可进一步简化。

同理，可以得出。
\[
\]
图画出了所示电源的对称相电压和线电压的相量图。则有
\[
\]
按照。

在对称情况下，三个相电流有效值相等。

\subsection{三角形联结}
图中，则三相总电压为
\[
\]
与上式对应的电压相量图。假如，则三相总电压为
\[
\]
相应的相量图。总电压有效值为相电压有效值的二倍。

当把负载接成三角形时，如果。

在三角形联结中。

在三角形联结中，用相量表示为
\[
\]

在负载端。

如果三个相电流是对称的，并由此得线电流与相电流的关系
\[
\]
可见在三角形联结中，当正弦相电流对称时，线电流也是对称的，其有效值等于，即
\[
\]
当按。

三角形联结中，在对称情况下可以记作
\[
\]

上面。

在三相电路。

一个三相负载应当接成。

\section{对称三相电路的计算}
图。

设，对负载中性点列电压方程
\[
\]
因为回路电压方程为
\[
\]
上式表明。

对于比较复杂。

\section{不对称三相电路示例}

\section{三相电路的功率}

\endinput